\documentclass[runningheads,a4paper]{llncs}

\usepackage[utf8]{inputenc}
 
\setcounter{tocdepth}{3}

\usepackage[english]{babel} 
\usepackage{graphicx}
\usepackage{grffile}
\usepackage{float}
\usepackage{multicol}
\usepackage{url}
 
\usepackage{titling}
\usepackage[hidelinks]{hyperref}
\setcounter{secnumdepth}{5}
%Margins
\usepackage[
margin=2cm,
includefoot
]{geometry}


\graphicspath{{img/}}

%Headers and Footers
\usepackage{fancyhdr}
\pagestyle{fancy}
\fancyhead{}
\fancyfoot{}
\fancyfoot[R]{\thepage}
\renewcommand{\headrulewidth}{0pt}
\renewcommand{\footrulewidth}{0pt}
 \setlength\parindent{24pt}
\begin{document}

%\title{NavUP
%	Software Requirements Specification}
%
%\author{Andries du Plooy (15226183), Vignesh Iyer (15031625), Thokozile Mabuzu (...), Mfana Masimula (12077713), Thabo Ntsoane  (15107532), Avinash Singh (14043778), Nicaedin Suklal (15207812)}
%
%\date{24 February 2017}
%
%\institute{University Of Pretoria - COS 301 Software Engineering}
%
%
%
%\authorrunning{Team Indigo}
%\titlerunning{SRC for NavUp}
%
	%Title Page
	\begin{titlepage}
		\begin{center}
			\includegraphics[width=10cm]{UP.jpg}  \\
			[1cm]
			\line(1,0){300} \\
			[0.3cm]
			\textsc{\Large
				NavUP\\
				Software Requirements Specification\\
			\hfill \break 24 February 2017
				%University of Pretoria
			}\\
			[0.1cm]
			\line(1,0){300} \\
			[0.7cm]
			\textsc{\Large
				Team Indigo
			} \\
			
			
			
		\end{center}
		
		\begin{center}
			\begin{multicols}{2}
				\textsc{\large\\
				Andries du Plooy\\ 
					15226183\\ 
				}
				
				\textsc{\large\\
				Vignesh Iyer\\
					 15031625\\ 
				}
				
				\textsc{\large\\
		        Mfana Masimula\\
				 	12077713\\ 
				}
				
				
				\textsc{\large\\
					Thabo Ntsoane\\
					15107532\\
				}
				
				\columnbreak
				
				\textsc{\large\\
					 Avinash Singh\\
					14043778\\
				}
				
				\textsc{\large\\
					Nicaedin Suklal\\
					15207812\\
				}

				
				\textsc{\large\\
					Thokozile Mabuzu\\
					???\\
				}
				
			\end{multicols}
			
			
			\textsc{	\\ \href{https://github.com/AndriesJacobus/SE-Indigo}{GitHub}
				\url{https://github.com/AndriesJacobus/SE-Indigo.git}}
			
		\end{center}
	\end{titlepage}
%\maketitle

\begingroup

\tableofcontents
\addcontentsline{toc}{section}{Table Of Contents}
\endgroup
\newpage


\section{Introduction}

\subsection{Purpose}
\paragraph*{The purpose of this Software Requirements Specification Document is to lay out the findings of a standardised requirements elicitation for the proposed application. In doing this, both the capabilities that the system is expected to deliver together, with the constraints on the solution space are described. The document will take into account requirements, as set out by both the developers and the prospective users, in order to lay a foundation for the up-coming phases of the Software Development Cycle. The various sections each provide details on specific types of requirements of the system. The intended audience include developers/peers(Team Indigo and other COS 301 teams), domain experts(COS 301 Lectures and the University of Pretoria Management), customers/end-users(students, staff and visitors at the University of Pretoria) who will perform a technical, expert  and customer review of this document respectively. }
\subsection{Scope}
\paragraph*{The name of the proposed software application is NavUP (NavigationUP). NavUP is intended to be an information service application that is to provide navigational information to visitors, employees and students of the University of Pretoria. This includes usual navigational functions (mapping and routing of campus), surveillance, user based information delivery and event calenders. NavUP will not extend its functionalities outside the Hatfield Campus and the University of Pretoria will not take any responsibility for any injury, damage or whatsoever that may be a result of the system. The application is used at your own risk. The objectives of the product include making the University of Pretoria feel more "like home" to all stakeholders alike. The application aims to facilitate a wide range of users who have had difficulties in the past with regard travelling within the premises. Ease of travel, accessibility of the various facilities, time optimisations, event/day-to-day coordination's, alleviating pedestrian traffic congestion problems, description of activities etc. are the joint objectives of the project. The goal is to integrate existing systems within the university like the Wi-Fi hotspots, IT-Services, Official University Website, Social Media and more.  Another goal is for the system to make use of human resources for its development and usage i.e. expertise from the COS 301, Honours groups of students and lectures to satisfy the former and make use of crowd-sourcing principles (using smart devices) to serve the latter. The benefits include smoothness of the day-to-day operations at the University which leads to greater productivity and efficiency among students and staff alike. It allows the university to increase its brand value in terms of the facilities it provides to keep external stakeholders pleased. National and Global Rankings of the university are bound to rise and consequently attract a larger group of people in the long run.}
\subsection{Definitions, Acronyms, and Abbreviations}
\paragraph{...}
\subsection{References}
\paragraph{...}
\subsection{Overview}
\paragraph*{The rest of this document will be solely focussed on elaborating on the requirements of the NavUP system. It is structured according to the IEEE STD 830-1998 standard. An overall description of the product will be followed by a specific detailed requirements description. Each item in both the sections represent the ideas as generated during compilation and there is possibilities of change during the implementation of the system.}

\section{Overall Description}

\subsection{Product Perspective}

\subsubsection{System Interfaces}
\paragraph{...}
\subsubsection{User Interfaces}
\paragraph{...}
\subsubsection{Hardware Interfaces}
\paragraph{...}
\subsubsection{Software Interfaces}
\paragraph{...}
\subsubsection{Communication Interfaces}
\paragraph{...}
\subsubsection{Memory}
\paragraph{...}
\subsubsection{Operations}
\paragraph{...}
\subsubsection{Site Adaptation Requirements}
\paragraph{...}

\subsection{Product Functions}
\paragraph{...}
\subsection{User Characteristics}
\paragraph{...}
\subsection{Constraints}
\paragraph{...}

% BEGIN Avinash
\subsection{Assumptions and Dependencies}
\let\labelitemi\labelitemii
The factors that affect the requirements are:
\begin{itemize}
	 
		\item Lack of connectivity -  Wi-Fi isn't available everywhere on campus (e.g. I.T building). thus assuming that Wi-Fi on campus will essentially be available everywhere in due time. \\
		
		\item Mobile devices - it is assumed that a user will have a mobile device that can run the system.
	 
\end{itemize}
% END Avinash


\section{Specific Requirements}

\subsection{External Interface Requirements}
\paragraph{...}
\subsection{Functional Requirements}
\paragraph{...}
\subsection{Performance Requirements}
\paragraph{...}
\subsection{Design Constraints}
\paragraph{...}

% BEGIN Avinash
 \subsection{Software System Attributes}
\subsubsection{Quality Requirements \\}
The following section will focus on the quality requirements of the overall system (mobile application) at its highest level of granularity. All the requirements defined, are for all components of the system.

\paragraph{Reliability \\}
The system should be reliable in the fact that the user should be able to reach his/her destination and that the information provided by the system is accurate and authentic. The system should also provide additional information aiding the reliability by giving a short description or history of the building or venue to make it more reliable. 

\paragraph{Security \\}
Although security is not a key feature, the system should not store any personal information or leak any information from the users device. The system should have access to the users Wi-Fi connections as well as cellular fall-back for enhanced navigation, and the system should not store any Wi-Fi information from the device or pull usernames or passwords from the users device. 

\paragraph{Availability \\}
 The system should be available to most recent types of user devices and operating systems on these devices. The system should always be on-line except if there is no Wi-Fi connectivity.

\paragraph{Flexibility \\}
 An important element of the system is its flexibility. 
 Each component of the system should be decoupled from every other component, thus allowing you to change certain components without affecting anything else. 
 For example it should be easy to add a new Lecture Venue to the system with little or no code modifications. 
 The main focus of the software is to allow the user to easily navigate to his/her destination with minimal effort.
 In spite of this the navigation route should be flexible in the sense there could be many different routes that the user can take and the system should be able to redirect or find another the best route that is in the users radius.

\paragraph{Maintainability \\}
 The system needs to be maintainable so that the future prospects of this system can expand to other campuses and in doing so code should be modular and well documented. Software and frameworks to be used should have a wide range of support and be established aiding the maintainability of the system.

\paragraph{Scalability \\}
 The system should be able to process data fast and data should not be lost so there should be a bus architecture to solve the data loss problem. The system should be coded to be scalable so that it can be deployed on many platforms such as Android, IOS and Microsoft with compatibility support.

\paragraph{Auditability \\}
 The system should make use of logs so that if there are any  queries or faults it can easily be traced or detected. These logs will also log navigational information, and for future prospects a user can opt in to share their experience where these log files could be sent to a remote server where the developers can analyse this data to improve the system. \\The following is necessary to be logged:
 \begin{itemize}
 	\item Program crashes
 	\item Navigation route (source and destination and some meta data [e.g. Time taken, estimated distance, etc.])
 \end{itemize}
 

\paragraph{Usability \\}
 The system should be easy to use and adhere to UI/UX standards with regards to the layout and user interfaces for interaction. The system should have easy navigation between functionality and views, as well as a intuitive way to display or simulate the user navigating to his/her destination.
 

% END Avinash


\subsection{Other Requirements}
\paragraph{...}

\section{Appendixes}
\paragraph{...}

\section{Index}
\paragraph{...}

\end{document}
